\chapter{Partitionnement (clustering), régimes et détection d'anomalies}\label{ch:ml:clustering-regimes-and-anomaly-detection}

Un tableau de bord de risque montre le marché dans son régime calme le vendredi précédant un krach. Le mardi, après réestimation, le même modèle montre
le krach commençant le vendredi. Les deux avaient raison sur ce qu'ils pouvaient savoir~: la première était une probabilité filtrée, calculée
à partir des données jusqu'au vendredi~; la seconde, une probabilité lissée, calculée en disposant du lundi suivant. Sur les régimes artificiels du chapitre, la
probabilité filtrée de stress la veille d'un changement de régime vaut en moyenne 0.03, pas plus que lors de n'importe quel jour calme~; la probabilité lissée
vaut en moyenne 0.35. Ce chapitre porte sur l'\omterm{def:ml:why-financial-machine-learning-is-different:learning}{apprentissage non supervisé}~: regrouper les actifs et les jours (\omterm{def:ml:clustering-regimes-and-anomaly-detection:cluster}{partitionnement}), décrire le marché
par des états cachés (modèles de régimes), et repérer ce qui ne cadre pas (\omterm{def:ml:clustering-regimes-and-anomaly-detection:anomaly}{détection d'anomalies}). Chacun est utile, et chacun appelle une
erreur particulière~: des groupes qui ne sont que du bruit, des régimes connus seulement a posteriori, des anomalies qui sont simplement rares.

\section{Partitionnement des actifs et des jours}

\begin{definition}[Partitionnement, k-moyennes, partitionnement hiérarchique]\label{def:ml:clustering-regimes-and-anomaly-detection:cluster}
Le \emph{partitionnement}\index{partitionnement} (clustering) répartit des objets en groupes dont les membres se ressemblent davantage entre eux qu'avec
les autres, sans étiquettes. \emph{k-moyennes}\index{k-moyennes} (k-means) choisit $k$ centres et affecte chaque objet au centre le plus proche, en alternant
les deux étapes pour minimiser la somme des carrés des distances intragroupes (algorithme de Lloyd). \emph{Partitionnement
hiérarchique}\index{partitionnement hiérarchique} (hierarchical clustering) fusionne pas à pas les groupes les plus proches en un arbre et le coupe au nombre souhaité
de groupes~; pour les actifs, la distance usuelle est $\sqrt{2(1-\rho_{ij})}$ entre séries de rendements (Mantegna, 1999), celle qui sous-tend
la parité de risque hiérarchique (HRP, livre~7, chapitre~26).
\end{definition}

Le chapitre regroupe les titres du marché synthétique de la firme (\texttt{firm.synthmkt} du livre~7) qui restent cotés pendant cinq
ans, dont les rendements dépendent d'un facteur de marché, dix facteurs sectoriels artificiels, des facteurs de style et un bruit de risque spécifique à queues épaisses. Le
facteur de marché est retiré en premier (le rendement de chaque titre moins son bêta multiplié par la moyenne à pondération égale), car il dominerait
toutes les corrélations. La \Cref{fig:ml:reg:clusters} évalue les \omterm{def:ml:clustering-regimes-and-anomaly-detection:cluster}{k-moyennes} et le \omterm{def:ml:clustering-regimes-and-anomaly-detection:cluster}{partitionnement hiérarchique} à liaison moyenne par l'indice de Rand ajusté
contre les secteurs artificiels (1 est parfait, 0 vaut le hasard), pour des fenêtres de trois mois à cinq ans et pour deux
intensités du facteur sectoriel~: la valeur par défaut, 8~\% par an, et 15~\%.

\begin{omfigure}
\begin{tikzpicture}
\begin{semilogxaxis}[width=11cm,height=6cm,xlabel={fenêtre (jours de bourse)},ylabel={indice de Rand ajusté},
  tick label style={font=\small},label style={font=\small},xmin=50,xmax=1500,ymin=-0.05,ymax=1.05,
  xtick={63,252,1260},xticklabels={63,252,1\,260},grid=major,grid style={gray!25},legend columns=2,legend cell align=left,
  legend style={font=\footnotesize,at={(0.5,-0.25)},anchor=north,draw=none,fill=none,
  /tikz/every even column/.append style={column sep=6pt}},table/col sep=comma]
\addplot[omThm,very thick,mark=*] table[x=window,y=km15]{figdata/ml/20-clustering-regimes-and-anomaly-detection/clusters.csv};
\addplot[omDef,very thick,mark=square*] table[x=window,y=hier15]{figdata/ml/20-clustering-regimes-and-anomaly-detection/clusters.csv};
\addplot[omThm!50,thick,dashed,mark=o,mark options={solid}] table[x=window,y=km08]{figdata/ml/20-clustering-regimes-and-anomaly-detection/clusters.csv};
\addplot[omDef!50,thick,dashed,mark=square,mark options={solid}] table[x=window,y=hier08]{figdata/ml/20-clustering-regimes-and-anomaly-detection/clusters.csv};
\legend{k-moyennes (facteur sectoriel 15~\%),partitionnement hiérarchique (facteur sectoriel 15~\%),k-moyennes (facteur sectoriel 8~\%),partitionnement hiérarchique (facteur sectoriel 8~\%)}
\end{semilogxaxis}
\end{tikzpicture}

\omcaption{Concordance des groupes fondés sur les rendements avec les secteurs artificiels, selon la longueur de la série de rendements et
l'intensité du facteur sectoriel. Données~: \texttt{ml\_regimes.clustering}.}\label{fig:ml:reg:clusters}
\end{omfigure}

Avec des secteurs marqués, le \omterm{def:ml:clustering-regimes-and-anomaly-detection:cluster}{partitionnement} les retrouve à partir de cinq ans de données (indice 0.99 et 1.00) et la moitié d'entre eux à
partir d'une année~; avec l'intensité par défaut, plus proche de ce que le risque spécifique laisse des effets sectoriels réels, il ne trouve presque
rien en un an (0.06) et de 0.15 à 0.21 en cinq ans. Le score de stabilité sur demi-échantillon mesure la même chose sans étiquettes~:
l'indice entre les regroupements des deux moitiés d'une année est de 0.32 à 0.35 avec des secteurs marqués et de 0.06 à 0.08 avec des secteurs faibles.
Un \omterm{def:ml:clustering-regimes-and-anomaly-detection:cluster}{partitionnement} qui ne résiste pas à la division de son propre échantillon décrit du bruit, et un desk qui s'en sert pour regrouper les risques ou
construire la parité de risque hiérarchique rééquilibre ses positions sur du bruit.

\section{Modèles de régimes}

\begin{definition}[Modèle de Markov caché, modèle à changement de régime, algorithme de Viterbi]\label{def:ml:clustering-regimes-and-anomaly-detection:hmm}
Un \emph{modèle de Markov caché}\index{modèle de Markov caché} (hidden Markov model, HMM) possède un état non observé qui suit une chaîne de Markov (livre~4,
chapitre~8) et des observations dont la loi dépend de l'état courant~; il est estimé par espérance-maximisation (livre~4,
chapitre~19), par l'algorithme de Baum--Welch (Rabiner, 1989). Un \emph{modèle à changement de régime}\index{modèle à changement de régime} (regime-switching model) est un HMM
pour des séries économiques, chaque état étant un régime doté de sa propre moyenne, de sa volatilité ou de sa dynamique (Hamilton, 1989~; Ang et Timmermann, 2012).
L'\emph{algorithme de Viterbi}\index{algorithme de Viterbi} (Viterbi algorithm) trouve le chemin d'états unique le plus vraisemblable compte tenu de l'ensemble des observations,
par programmation dynamique.
\end{definition}

L'indice du chapitre bascule entre un régime calme (moyenne quotidienne de 5 points de base, volatilité de 0.8~\%) et un régime de stress (moyenne
$-10$ points de base, volatilité de 2.5~\%), en restant dans chacun avec une probabilité de 0.99 et de 0.95 par jour~; vingt ans sont simulés, et
le régime de stress occupe un cinquième du temps. Un HMM gaussien à deux états estimé sur les dix premières années identifie le calme avec une volatilité de 0.80~\% et
le régime de stress avec 2.60~\%, des probabilités de persistance de 0.988 et 0.958. Un modèle à trois états ajoute 5.7 à la log-vraisemblance avec six paramètres
de plus, moins que la pénalité de 23.5 du critère d'information bayésien, et consacre son troisième état à scinder le régime de stress en deux
(2.28~\% et 2.65~\%).

\section{Filtrage contre lissage~: ce que sait un modèle de régimes en temps réel}

Un HMM fournit trois réponses sur l'état à la date $t$. La probabilité \emph{filtrée} utilise les données jusqu'à $t$ (la récursion avant
du \cref{lst:ml:reg:hmm})~; la probabilité \emph{lissée} utilise toutes les données, y compris celles qui suivent $t$ (la
passe arrière)~; le chemin de Viterbi est la meilleure séquence compte tenu de l'ensemble des données. Seule la première était disponible en $t$. Sur les dix années de test,
l'estimation filtrée (régime de stress lorsque sa probabilité dépasse un demi) est juste pour 95.9~\% des jours, la lissée pour 97.7~\% et le
chemin de Viterbi pour 97.5~\%~; une règle fondée sur la volatilité glissante de 20 jours est juste pour 89.8~\%. Après chacun des 20 basculements vers le
régime de stress, le filtre l'indique après une médiane d'un jour (moyenne 1.5, un épisode bref manqué), l'estimation lissée le jour même
(moyenne 0.8). La \Cref{fig:ml:reg:switch} montre un basculement~: la probabilité lissée augmente avant le changement de régime, car les
jours qui suivent le basculement font ressembler au stress les jours qui le précèdent.

\begin{omfigure}
\begin{tikzpicture}
\begin{axis}[width=12cm,height=6cm,xlabel={jours depuis un basculement vers le stress},ylabel={probabilité de stress},
  tick label style={font=\small},label style={font=\small},xmin=-30,xmax=49,ymin=-0.02,ymax=1.05,grid=major,
  grid style={gray!25},legend columns=3,legend cell align=left,legend style={font=\footnotesize,at={(0.5,-0.25)},anchor=north,
  draw=none,fill=none,/tikz/every even column/.append style={column sep=6pt}},table/col sep=comma]
\addplot[gray!60,fill=gray!15,const plot,draw=none] table[x=day,y=state]{figdata/ml/20-clustering-regimes-and-anomaly-detection/switch.csv} \closedcycle;
\addplot[omThm,very thick] table[x=day,y=filtered]{figdata/ml/20-clustering-regimes-and-anomaly-detection/switch.csv};
\addplot[omDef,very thick,dashed] table[x=day,y=smoothed]{figdata/ml/20-clustering-regimes-and-anomaly-detection/switch.csv};
\legend{régime de stress réel,filtrée (connue ce jour-là),lissée (connue à la fin)}
\end{axis}
\end{tikzpicture}

\omcaption{Probabilités filtrée et lissée de stress autour d'un basculement dans les années de test (le premier basculement que le filtre a mis
au moins deux jours à le repérer). Données~: \texttt{ml\_regimes} (\texttt{fig\_regimes.py}).}\label{fig:ml:reg:switch}
\end{omfigure}

L'écart est un biais d'anticipation (look-ahead bias, livre~7, chapitre~3) que nul horodatage ne révèle. La veille d'un basculement, la probabilité filtrée
de stress vaut en moyenne 0.032, inférieure à sa moyenne de 0.041 sur l'ensemble des jours calmes~: le filtre ne voit pas venir le basculement. La
probabilité lissée vaut en moyenne 0.353 ces jours-là, contre 0.019 les jours calmes. Un backtest qui ajuste l'exposition pour viser une
volatilité cible à partir des probabilités de régime affiche un ratio de Sharpe de 1.00 avec la probabilité filtrée de la veille, reportée
par la matrice de transition, et de 1.21 avec la probabilité lissée du jour même~; la stratégie d'achat et de conservation affiche 0.91. Un modèle de régimes
utilisé pour le trading ou le risque doit être exécuté comme un filtre, réestimé uniquement sur les données disponibles à chaque date, et ses statistiques
(la volatilité d'un régime, sa durée moyenne) doivent être calculées à partir des probabilités filtrées.

\section{Détection d'anomalies}

\begin{definition}[Détection d'anomalies, forêt d'isolement, courbe précision--rappel]\label{def:ml:clustering-regimes-and-anomaly-detection:anomaly}
Une \emph{détection d'anomalies}\index{détection d'anomalies} (anomaly detection) attribue un score aux observations selon leur écart au gros des données, sans
étiquettes pour les anomalies. Une \emph{forêt d'isolement}\index{forêt d'isolement} (isolation forest) construit des arbres aléatoires qui coupent sur des variables tirées au hasard
à des valeurs tirées au hasard~; les points isolés après peu de coupures sont anormaux (Liu, Ting et Zhou, 2008). La
\emph{courbe précision--rappel}\index{courbe précision--rappel} (precision--recall curve) trace, à mesure que le seuil d'alerte varie, la part des alertes qui sont
vraies (précision) contre la part des cas vrais qui déclenchent une alerte (rappel)~; lorsque les positifs sont rares, elle est plus informative
que la courbe ROC (Davis et Goadrich, 2006).
\end{definition}

Le livre~9 (chapitre~29) a construit des détecteurs de spoofing sur des journées de compte~: 100 épisodes semés parmi 11\,500 journées de compte légitimes,
dont des teneurs de marché qui annulent presque tout et des fournisseurs de liquidité du carnet profond dont les gros ordres s'exécutent rarement.
Le \Cref{tab:ml:reg:anomaly} ajoute deux scores non supervisés sur cinq variables des mêmes journées de compte (les logarithmes des nombres de
petits et de gros ordres, leurs taux d'exécution, et la part des grosses annulations qui suivent une exécution du côté opposé)~: une
\omterm{def:ml:clustering-regimes-and-anomaly-detection:anomaly}{forêt d'isolement} (\cref{lst:ml:reg:iforest}) et l'erreur de reconstruction d'un petit \omterm{def:ml:cross-sectional-deep-models:ae}{autoencodeur} (chapitre~9).

\begin{table}[htbp]
\centering\small
\begin{tabular}{@{}lcc@{}}
\toprule
détecteur & rappel à 1~\% de faux positifs & précision\\
\midrule
ratio ordres sur transactions (règle) & 0.00 & 0.00\\
écart de taux d'exécution (règle) & 0.37 & 0.24\\
annulations après une exécution (règle) & 0.51 & 0.31\\
écart $\times$ annulations (règle) & 0.77 & 0.40\\
\omterm{def:ml:clustering-regimes-and-anomaly-detection:anomaly}{forêt d'isolement} & 1.00 & 0.47\\
\omterm{def:ml:cross-sectional-deep-models:ae}{autoencodeur} & 0.32 & 0.22\\
\bottomrule
\end{tabular}
\caption{Détecteurs de spoofing sur 11\,600 journées de compte (100 épisodes semés), au seuil qui signale 1~\% des journées de compte
légitimes. Données~: \texttt{ml\_regimes.anomalies}.}\label{tab:ml:reg:anomaly}
\end{table}

\section{Détecteurs de surveillance}

La \omterm{def:ml:clustering-regimes-and-anomaly-detection:anomaly}{forêt d'isolement} retrouve chaque épisode semé au taux de 1~\% de faux positifs, et sa précision, 0.47, est le maximum que tout
détecteur peut atteindre à ce niveau~: 115 journées de compte légitimes sont signalées quel que soit le détecteur, contre au plus 100 vraies. L'\omterm{def:ml:cross-sectional-deep-models:ae}{autoencodeur}
fait moins bien que les règles. Trois précautions relativisent ce résultat. La forêt a reçu les variables choisies par les
auteurs des règles, qui codent ce à quoi ressemble le spoofing~; sur des journaux d'ordres bruts, elle repérerait l'inhabituel, et non
la manipulation. Les spoofeurs semés forment une population réduite et distincte~; lorsqu'un schéma manipulateur est courant, ou que
les traders légitimes qui lui ressemblent sont nombreux, un score non supervisé les classe ensemble. Enfin, la précision à taux de
faux positifs fixé est ce avec quoi vit une équipe de conformité~: un détecteur au rappel parfait lui remet tout de même davantage de
journées de compte innocentes que coupables. Les scores non supervisés gagnent leur place comme second filet sous les règles, revus par des personnes, et comme moyen de repérer des schémas qu'aucune règle ne nomme encore.

\begin{method}[Apprentissage non supervisé sur les données de marché]\label{met:ml:reg:method}
\begin{enumerate}
\item Retirer le facteur commun avant de partitionner les rendements, et indiquer un score de stabilité sur les sous-échantillons pour chaque
\omterm{def:ml:clustering-regimes-and-anomaly-detection:cluster}{partitionnement}.
\item Estimer les modèles de régimes sur les données disponibles à chaque date et n'utiliser que les probabilités filtrées~; présenter les résultats lissés comme
des descriptions de l'histoire, jamais comme des backtests.
\item Choisir le nombre de régimes par une vraisemblance hors échantillon (out-of-sample) ou un critère d'information, et vérifier que les régimes diffèrent
en quelque chose sur quoi le desk agit.
\item Évaluer les scores d'anomalies comme des détecteurs~: rappel et précision au taux de faux positifs que les relecteurs
peuvent absorber, face aux règles, sur des données comportant des cas semés et des cas qui leur ressemblent.
\end{enumerate}
\end{method}

\section{Tutoriel~: dans quel régime sommes-nous~?}\label{tut:ml:clustering-regimes-and-anomaly-detection:tut}

\begin{tutorial}
\textbf{Objectif.} Partitionner les noms synthétiques, estimer et exécuter un modèle de régimes en filtre et en lisseur, et évaluer les détecteurs non supervisés
contre les règles de surveillance. \textbf{État final~:} les \cref{fig:ml:reg:clusters,fig:ml:reg:switch},
et le \cref{tab:ml:reg:anomaly}.
\begin{tutsteps}
\item \textbf{\omterm{def:ml:clustering-regimes-and-anomaly-detection:cluster}{Partitionnement hiérarchique} sur la distance de corrélation.}
\omcode{code/firm/regimes/firm_regimes.py}{43}{50}{Liaison moyenne sur la distance de Mantegna.}\label{lst:ml:reg:hier}
\item \textbf{Filtrage et lissage.}
\omcode{code/firm/regimes/firm_regimes.py}{84}{109}{La récursion avant (filtre) et la passe arrière (lisseur).}\label{lst:ml:reg:hmm}
\item \textbf{Une \omterm{def:ml:clustering-regimes-and-anomaly-detection:anomaly}{forêt d'isolement}.}
\omcode{code/firm/regimes/firm_regimes.py}{161}{166}{Scores d'anomalies de la forêt d'isolement.}\label{lst:ml:reg:iforest}
\item \textbf{Exécuter} \texttt{ml\_regimes.clustering()}, \texttt{regimes()}, \texttt{vol\_targeting()}, \texttt{anomalies()} et
\texttt{fig\_regimes.py}.
\end{tutsteps}
\textbf{Pour aller plus loin.} Réestimer le HMM chaque année sur une fenêtre croissante et filtrer l'année suivante~; doubler le nombre de
fournisseurs de liquidité du carnet profond légitimes et relancer les scores d'anomalies.
\end{tutorial}

\section{Construction~: partitionnement, régimes et scores d'anomalies}\label{bld:ml:clustering-regimes-and-anomaly-detection:regimes}

\begin{build}
\bfield{Rôle} Outils non supervisés dont les sorties sont connues en temps réel et évalués comme des détecteurs.
\bfield{Interface} \texttt{market\_residuals}, \texttt{corr\_distance}, \texttt{kmeans\_clusters}, \texttt{hier\_clusters},
\texttt{ari}, \texttt{stability}~; \texttt{regime\_series}~; \texttt{GaussianHMM(k)} avec \texttt{fit}, \texttt{filter},
\texttt{smooth}, \texttt{viterbi}~; \texttt{iforest\_scores}, \texttt{autoencoder\_scores}~; avec la
\texttt{firm.surveil.tpr\_at\_fpr} du livre~9.
\bfield{Règles} Probabilités filtrées uniquement dans tout ce qui opère sur le marché ou rapporte un risque~; chaque \omterm{def:ml:clustering-regimes-and-anomaly-detection:cluster}{partitionnement} avec son score de stabilité~;
des scores d'anomalies avec rappel et précision à un taux de faux positifs fixé.
\bfield{Tests d'acceptation} \texttt{code/firm/regimes/tests/}~: le \omterm{def:ml:clustering-regimes-and-anomaly-detection:cluster}{partitionnement} retrouve des groupes semés bien séparés~; le HMM retrouve les paramètres
semés et son filtre n'utilise aucune donnée future (il est inchangé lorsque les données ultérieures changent)~; le lissage et le filtrage concordent
le dernier jour~; l'\omterm{def:ml:clustering-regimes-and-anomaly-detection:hmm}{algorithme de Viterbi} coïncide avec la force brute sur une courte série~; la \omterm{def:ml:clustering-regimes-and-anomaly-detection:anomaly}{forêt d'isolement} classe en tête les valeurs aberrantes semées.
\bfield{Pour aller plus loin} Régimes dotés de leur propre dynamique (autorégressifs ou GARCH au sein des états)~; réestimation en ligne~; \omterm{def:ml:clustering-regimes-and-anomaly-detection:cluster}{partitionnement} des journées
plutôt que des actifs.
\end{build}

\begin{omsources}
\item J.~D.~Hamilton, ``A new approach to the economic analysis of nonstationary time series and the business cycle'',
\emph{Econometrica} 57(2), 1989.
\item A.~Ang et A.~Timmermann, ``Regime changes and financial markets'', \emph{Annual Review of Financial Economics} 4, 2012.
\item L.~R.~Rabiner, ``A tutorial on hidden Markov models and selected applications in speech recognition'',
\emph{Proceedings of the IEEE} 77(2), 1989.
\item F.~T.~Liu, K.~M.~Ting et Z.-H.~Zhou, ``Isolation forest'', \emph{IEEE International Conference on Data Mining}, 2008.
\item S.~Lloyd, ``Least squares quantization in PCM'', \emph{IEEE Transactions on Information Theory} 28(2), 1982.
\item R.~N.~Mantegna, ``Hierarchical structure in financial markets'', \emph{European Physical Journal B} 11, 1999.
\item J.~Davis et M.~Goadrich, ``The relationship between precision-recall and ROC curves'', \emph{ICML}, 2006.
\end{omsources}

\section{Exercices}

\begin{exercise}[$\star$]\label{exo:ml:clustering-regimes-and-anomaly-detection:1}
Deux noms ont une corrélation des rendements de 0.5. Calculer leur distance de Mantegna, puis déterminer la distance
des noms parfaitement corrélés et celle des noms non corrélés.
\end{exercise}

\begin{exercise}[$\star$]\label{exo:ml:clustering-regimes-and-anomaly-detection:2}
Avec des probabilités de rester de 0.988 et 0.958, quelles sont les durées moyennes des régimes calme et de stress, et quelle est la
part de temps passée en stress à long terme~?
\end{exercise}

\begin{exercise}[$\star$]\label{exo:ml:clustering-regimes-and-anomaly-detection:3}
Avec un taux de faux positifs de 1~\% sur 11\,500 jours de compte légitimes, quelle est la précision maximale qu'un détecteur quelconque peut atteindre avec 100
cas réels~?
\end{exercise}

\begin{exercise}[$\star\star$]\label{exo:ml:clustering-regimes-and-anomaly-detection:4}
Pourquoi la probabilité filtrée de stress la veille d'un changement de régime n'est-elle pas plus élevée que lors d'une journée calme moyenne~?
\end{exercise}

\begin{exercise}[$\star\star$]\label{exo:ml:clustering-regimes-and-anomaly-detection:5}
Pourquoi faut-il retirer le facteur de marché avant de partitionner les rendements~?
\end{exercise}

\begin{exercise}[$\star\star$]\label{exo:ml:clustering-regimes-and-anomaly-detection:6}
\emph{Trouver la faille.} «~La probabilité de stress de notre modèle à deux régimes, calibrée sur 2000--2024, s'est manifestée quelques semaines avant chacun des
trois grands drawdowns, elle constitue donc un signal d'alerte précoce.~»
\end{exercise}

\begin{exercise}[$\star\star\star$]\label{exo:ml:clustering-regimes-and-anomaly-detection:7}
\emph{Programmation.} Comparer le modèle à trois états et le modèle à deux états sur les années de test, au moyen de la log-vraisemblance des données de test
sous chaque modèle calibré (en effectuant le filtrage avant à partir des paramètres calibrés). Lequel l'emporte hors échantillon~?
\end{exercise}

\begin{exercise}[$\star\star\star$]\label{exo:ml:clustering-regimes-and-anomaly-detection:8}
Établir la récursion avant $\alpha_t(j) = \big(\sum_i\alpha_{t-1}(i)P_{ij}\big)b_j(x_t)$ pour la probabilité conjointe des
observations jusqu'à $t$ et de l'état en $t$, puis montrer que sa normalisation donne la probabilité filtrée.
\end{exercise}

\section{Problème~: Dans quel régime sommes-nous~?}

\begin{problem}[{Problème du week-end --- filtré, lissé et rare}]\label{pb:ml:clustering-regimes-and-anomaly-detection:1}
Les partitions, les régimes et les détecteurs du chapitre.

\medskip\noindent\textbf{Partie I --- Partitions.}
\begin{enumerate}
\item Dans quelle mesure les partitions retrouvent-elles les secteurs, selon la fenêtre et la force du signal~?
\item Que dit le score de stabilité, et pourquoi est-ce celui qu'il faut rapporter sur des données réelles~?
\item À quoi servent les partitions sur un desk, et à quoi ne servent-elles pas~?
\item Quelle méthode de \omterm{def:ml:clustering-regimes-and-anomaly-detection:cluster}{partitionnement} choisirait-on, et pourquoi~?
\end{enumerate}

\medskip\noindent\textbf{Partie II --- Régimes.}
\begin{enumerate}[resume]
\item Que trouve le modèle à deux états calibré, en comparaison des régimes simulés~?
\item Pourquoi ne pas retenir trois états~?
\item Quelle est l'exactitude de chacune des quatre estimations de régime sur les années de test~?
\item Quels sont les délais avant les passages en stress~?
\end{enumerate}

\medskip\noindent\textbf{Partie III --- Avec le recul.}
\begin{enumerate}[resume]
\item Que disent les probabilités filtrées et lissées la veille d'un changement de régime~?
\item Dans quelle mesure la probabilité lissée embellit-elle un backtest de ciblage de volatilité~?
\item Comment faut-il calculer les statistiques par régime~?
\item Où le lissage se cache-t-il ailleurs dans un pipeline de recherche~?
\end{enumerate}

\medskip\noindent\textbf{Partie IV --- Le verdict.}
\begin{enumerate}[resume]
\item Énoncer le \emph{résultat nommé}~: l'exactitude et le délai de détection du régime filtré par rapport au régime lissé, ainsi que la précision des détecteurs
d'anomalies à un taux de faux positifs fixé, par rapport aux détecteurs à règles.
\item À quoi sert un modèle de régime, si ce n'est à prédire~?
\item Laisserait-on la \omterm{def:ml:clustering-regimes-and-anomaly-detection:anomaly}{forêt d'isolement} remplacer les règles~?
\item Comment examinerait-on les alertes de la \omterm{def:ml:clustering-regimes-and-anomaly-detection:anomaly}{forêt d'isolement}~?
\item Que surveillerait-on pour savoir qu'un modèle de régime a cessé de fonctionner~?
\item Comment les partitions alimentent-elles la parité de risque hiérarchique, et que se passe-t-il lorsqu'elles ne sont que du bruit~?
\item Quelles étiquettes transformeraient le problème d'anomalies en problème supervisé, et à quel coût~?
\item En une phrase~: que sait un modèle d'aujourd'hui~?
\end{enumerate}
\end{problem}

\section{Questions d'entretien}

\begin{interviewq}[$\star$ \iqroles{chercheur}]\label{iq:ml:clustering-regimes-and-anomaly-detection:1}
Quelle est la différence entre les probabilités d'états filtrées et lissées dans un HMM~?
\end{interviewq}

\begin{interviewq}[$\star\star$ \iqroles{chercheur, ingénieur ML}]\label{iq:ml:clustering-regimes-and-anomaly-detection:2}
Comment partitionneriez-vous 500 actions selon leurs rendements, et comment sauriez-vous que les partitions sont réelles~?
\end{interviewq}

\begin{interviewq}[$\star\star$ \iqroles{chercheur}]\label{iq:ml:clustering-regimes-and-anomaly-detection:3}
Expliquez en quelques phrases l'algorithme de Baum--Welch.
\end{interviewq}

\begin{interviewq}[$\star\star$ \iqroles{ingénieur ML}]\label{iq:ml:clustering-regimes-and-anomaly-detection:4}
Comment fonctionne une \omterm{def:ml:clustering-regimes-and-anomaly-detection:anomaly}{forêt d'isolement}, et quels sont ses angles morts~?
\end{interviewq}

\begin{interviewq}[$\star\star$ \iqroles{chercheur, trader}]\label{iq:ml:clustering-regimes-and-anomaly-detection:5}
Une surcouche à changement de régime a fait passer le ratio de Sharpe backtesté d'une stratégie de 0.9 à 1.2. Que vérifiez-vous~?
\end{interviewq}

\begin{interviewq}[$\star\star\star$ \iqroles{ingénieur ML}]\label{iq:ml:clustering-regimes-and-anomaly-detection:6}
Concevez un système de \omterm{def:ml:clustering-regimes-and-anomaly-detection:anomaly}{détection d'anomalies} pour la surveillance des transactions. Comment l'évaluez-vous sans disposer de beaucoup d'étiquettes~?
\end{interviewq}
